\documentclass[11pt]{article}

\usepackage[margin=1in]{geometry}
\usepackage{amsmath}
\usepackage{amssymb}
\usepackage{siunitx}
\usepackage{booktabs}
\usepackage[hidelinks]{hyperref}

\title{Low Cycle Fatigue: Core Physics and Equations}
\author{Low Cycle Fatigue MCP Project}
\date{\today}

\begin{document}
\maketitle

\begin{abstract}
A brief reference for the physics and governing equations behind low cycle fatigue (LCF)
analysis. This document defines the quantities the analysis tool computes. All analysis uses
\emph{true} stress and \emph{true} strain.
\end{abstract}

\section{Background}

Fatigue is the progressive, localized damage a material accumulates under cyclic loading,
leading to failure at stresses below the monotonic ultimate strength. \emph{Low cycle fatigue}
(LCF) refers to the regime in which each cycle imposes significant plastic deformation, so
failure occurs in a relatively small number of cycles (typically $10^{1}$--$10^{4}$). LCF is
strain-controlled: a specimen is cycled between fixed strain limits, and the stress response
is measured.

Throughout, we use true stress $\sigma$ and true strain $\varepsilon$, related to engineering
quantities by
\begin{equation}
  \varepsilon = \ln(1 + \varepsilon_\text{eng}),
  \qquad
  \sigma = \sigma_\text{eng}\,(1 + \varepsilon_\text{eng}).
\end{equation}

\section{Cyclic Stress--Strain Response}

LCF tests are typically run under \emph{fully reversed} strain control, characterized by the
strain ratio
\begin{equation}
  R = \frac{\varepsilon_{\min}}{\varepsilon_{\max}} = -1,
\end{equation}
i.e.\ each cycle swings symmetrically between equal tensile and compressive strains. Under this
loading the stress--strain path traces a closed \emph{hysteresis loop} each cycle. The total
strain amplitude decomposes into elastic and plastic parts:
\begin{equation}
  \frac{\Delta\varepsilon_t}{2}
  = \frac{\Delta\varepsilon_e}{2} + \frac{\Delta\varepsilon_p}{2},
  \qquad
  \frac{\Delta\varepsilon_e}{2} = \frac{\Delta\sigma}{2E},
\end{equation}
where $\Delta\sigma$ is the stress range, $E$ is Young's modulus, and the plastic amplitude
$\Delta\varepsilon_p/2$ is measured from the width of the loop.

As cycling proceeds the material typically \emph{hardens} (peak stress rises) and then
\emph{softens} (peak stress falls) toward failure. A tension/compression asymmetry in the
peak stresses indicates direction-dependent deformation, quantified by the ratio
\begin{equation}
  R_{TC} = \frac{\lvert \sigma_{\max,\text{tension}} \rvert}
                {\lvert \sigma_{\max,\text{compression}} \rvert}.
\end{equation}
Note that even under fully reversed strain control ($R=-1$), this stress ratio need not equal
unity; $R_{TC}\neq 1$ reflects a nonzero mean stress in the cycle (see Morrow's correction,
\S\ref{sec:morrow}).

\section{Cyclic Energy Density}

The area enclosed by one hysteresis loop is the mechanical energy dissipated per unit volume
per cycle:
\begin{equation}
  W = \oint \sigma(\varepsilon)\,\mathrm{d}\varepsilon
  \qquad \left[\si{\joule\per\cubic\meter}\right].
\end{equation}
It is reported at the peak-hardened state and at the half-life (stable) cycle, conventionally
in \si{\mega\joule\per\cubic\meter}.

\section{Strain--Life Relations}

\subsection{Total strain--life}
Fatigue life is modeled as the sum of elastic (Basquin) and plastic (Coffin--Manson)
contributions versus the reversals to failure $2N_f$:
\begin{equation}
  \frac{\Delta\varepsilon_t}{2}
  = \underbrace{\frac{\sigma'_f}{E}\,(2N_f)^{b}}_{\text{elastic (Basquin)}}
  + \underbrace{\varepsilon'_f\,(2N_f)^{c}}_{\text{plastic (Coffin--Manson)}}.
\end{equation}

\subsection{Basquin (elastic branch)}
\begin{equation}
  \frac{\Delta\sigma}{2} = \sigma'_f\,(2N_f)^{b},
\end{equation}
with fatigue strength coefficient $\sigma'_f$ and exponent $b$ obtained from a linear fit on
$\log(\Delta\sigma/2)$ versus $\log(2N_f)$.

\subsection{Coffin--Manson (plastic branch)}
\begin{equation}
  \frac{\Delta\varepsilon_p}{2} = \varepsilon'_f\,(2N_f)^{c},
\end{equation}
with fatigue ductility coefficient $\varepsilon'_f$ and exponent $c$ obtained from a linear fit
on $\log(\Delta\varepsilon_p/2)$ versus $\log(2N_f)$.

\subsection{Morrow mean-stress correction}
\label{sec:morrow}
The base strain--life relation assumes a zero mean stress. When the cycle carries a nonzero
mean stress $\sigma_m$ (e.g.\ from the tension/compression asymmetry $R_{TC}\neq 1$, even at a
strain ratio of $R=-1$), Morrow's relation corrects the elastic (Basquin) term by shifting the
fatigue strength coefficient by $\sigma_m$:
\begin{equation}
  \frac{\Delta\varepsilon_t}{2}
  = \frac{\sigma'_f - \sigma_m}{E}\,(2N_f)^{b}
  + \varepsilon'_f\,(2N_f)^{c},
\end{equation}
where the mean stress of a cycle is
\begin{equation}
  \sigma_m = \frac{\sigma_{\max} + \sigma_{\min}}{2}.
\end{equation}
A tensile mean stress ($\sigma_m > 0$) reduces fatigue life, while a compressive mean stress
($\sigma_m < 0$) extends it. The relation reduces to the base strain--life equation when
$\sigma_m = 0$. Morrow's correction leaves the plastic (Coffin--Manson) term unchanged, on the
assumption that mean stress relaxes in the high-plasticity regime and chiefly affects the
elastic, longer-life contribution.

\section{Cyclic Stress--Strain Curve (Ramberg--Osgood)}

The stabilized stress--strain behavior across amplitudes is described by
\begin{equation}
  \Delta\varepsilon_t
  = \frac{\Delta\sigma}{E}
  + \left(\frac{\Delta\sigma}{K'}\right)^{1/n'},
\end{equation}
where the plastic term is fit in linearized form
\begin{equation}
  \frac{\Delta\sigma}{2} = K'\left(\frac{\Delta\varepsilon_p}{2}\right)^{n'}
  \quad\Longleftrightarrow\quad
  \log\!\left(\frac{\Delta\sigma}{2}\right)
  = n'\log\!\left(\frac{\Delta\varepsilon_p}{2}\right) + \log K',
\end{equation}
giving the cyclic strength coefficient $K'$ and cyclic strain-hardening exponent $n'$.

\section{Consistency Relations}

When the fitted constants are mutually consistent, the following hold and serve as validity
checks:
\begin{equation}
  n' \approx \frac{b}{c},
  \qquad
  K' \approx \frac{\sigma'_f}{(\varepsilon'_f)^{\,b/c}}.
\end{equation}

\section{Symbol Reference}

\begin{table}[h]
\centering
\begin{tabular}{lll}
\toprule
Symbol & Quantity & Unit \\
\midrule
$\sigma,\ \varepsilon$ & true stress, true strain & \si{\pascal}, --- \\
$\Delta\varepsilon_t/2$ & total strain amplitude & --- \\
$\Delta\varepsilon_e/2$ & elastic strain amplitude & --- \\
$\Delta\varepsilon_p/2$ & plastic strain amplitude & --- \\
$\Delta\sigma/2$ & stress amplitude & \si{\pascal} \\
$\sigma_m$ & mean stress (Morrow) & \si{\pascal} \\
$R$ & strain ratio ($\varepsilon_{\min}/\varepsilon_{\max}$, $-1$ fully reversed) & --- \\
$R_{TC}$ & tension/compression peak-stress ratio & --- \\
$2N_f$ & reversals to failure & --- \\
$E$ & Young's modulus & \si{\pascal} \\
$W$ & cyclic energy density & \si{\mega\joule\per\cubic\meter} \\
$\sigma'_f,\ b$ & fatigue strength coefficient, exponent & \si{\pascal}, --- \\
$\varepsilon'_f,\ c$ & fatigue ductility coefficient, exponent & ---, --- \\
$K',\ n'$ & cyclic strength coefficient, hardening exponent & \si{\pascal}, --- \\
\bottomrule
\end{tabular}
\end{table}

\end{document}
